\section{Proof of the primitive sampling theorem}
\label{samp:proofs}

This appendix proves Theorem~\ref{samp:staircase}.  The main text gave
its probabilistic mechanism; here we justify the limiting distribution,
the exact Lebesgue decomposition, and the dimension statement.

\begin{proof}
We first obtain the limiting probability model.  Take a pair of
Haar-uniform profinite integers $(U,V)$, conditioned on having no
common prime divisor.  This event has probability
$\prod_p(1-p^{-2})=6/\pi^2$.  At each prime $p$, the residues $(U,V)$
are therefore uniform among the $p^2-1$ pairs other than $(0,0)$;
these choices are independent at distinct primes.  Independently
choose fair bits $B_0,B_1,\ldots$.  The candidate limit is the law of
\begin{equation}
 Z_g=\sum_{j\ge0}\bigl(1-h_g(U+jV)\bigr)B_j2^{-j-1}.
 \label{samp:model}
\end{equation}
The least prime of each $U+jV$ exists almost surely.  For example,
under this local law
$\Prob(p\mid U+jV)=1/(p+1)$, and independence over primes shows
that the probability of avoiding every prime is zero.

For completeness, the analytic input that produces the fair bits
is the following variable-step cancellation: for any nonempty finite
$J\subset\{0,1,\ldots\}$ and any fixed function $P(a,d)$ periodic
in both variables,
\begin{equation}
 \frac1{N^2}\sum_{a,d\le N}P(a,d)
       \prod_{j\in J}\lambda(a+jd)\longrightarrow0.
 \label{samp:cancellation}
\end{equation}
Theorem~2.5 of Frantzikinakis--Host \cite{sampFrantzikinakisHost2017}
gives vanishing Gowers norms of every fixed order for the aperiodic
function $\lambda$; their Lemma~9.6 controls products on independent
linear forms.  To include $P$, expand it in its finite Fourier
series.  For $j_0\in J$ a phase $e((ra+sd)/Q)$ factors as
$e(r(a+j_0d)/Q)e((s-rj_0)d/Q)$, where $e(t)=e^{2\pi it}$.
Absorb the first factor into $\lambda(a+j_0d)$ and the second into
the additional form $d$.  Multiplication by a linear phase preserves
Gowers norms of order at least two.  Extend the first function evenly
from positive integers, as required by that lemma.  The forms
$a+jd$ and $d$ are pairwise independent, proving
\eqref{samp:cancellation}.  Thus no assertion about correlations
at a fixed step is being used.

Fix a prefix length $k$.  Truncate both the mask and the coprimality
condition at primes at most $R$.  All mask coefficients and
restrictions are now periodic in $(a,d)$.  Expand a prefix indicator
using $1-2x_g=h_g+(1-h_g)\lambda$.  Every term containing a
Liouville factor vanishes by \eqref{samp:cancellation}; the other
terms converge by the Chinese remainder theorem.  These are exactly
the prefix probabilities in \eqref{samp:model}, with the same
truncations.  The proportion of pairs missed by truncated
coprimality is at most $\sum_{p>R}p^{-2}$, since
$\lfloor N/p\rfloor^2\le N^2/p^2$.  A mask error at $a+jd$
requires that integer to avoid all primes at most $R$.  For fixed
$R$, its limiting proportion in the square is
$\prod_{p\le R}(1-1/p)$, which tends to zero.  Summing this error
over the $k$ positions and then removing both truncations proves
prefix convergence.  A continuous function of the sampled real is
uniformly approximated by a function of its first $k$ digits, since
the remaining binary tail is at most $2^{-k}$.  This proves weak
convergence to the law of $Z_g$.

Next fix a primitive profinite pair $(U,V)$.  A lane $p$ is visited
by $U+jV$, $j\ge0$, if and only if $p\nmid V$.  If $p\mid V$,
primitivity gives $p\nmid U$, so the lane is avoided.  Otherwise
prescribe the unique root modulo $p$ and avoid the roots modulo
each smaller prime $q$ with $q\nmid V$.  There is an allowed residue
at each such $q$; primes dividing $V$ already divide no term.
The Chinese remainder theorem gives an index with least prime $p$,
and the same event repeats periodically in $j$.

All digits of the progression are open precisely when $p\mid V$
for every $p\in\mathcal F$.  Locally this has probability
$(p-1)/(p^2-1)=1/(p+1)$, proving that the all-open event has
probability $L_g$.  On it, $Z_g$ is uniform on $[0,1]$.
An all-closed progression would require $p\mid V$ for every
successful prime.  There are infinitely many such primes by
\eqref{samp:success-hypothesis}, so that event has probability zero.
Almost every fiber therefore has an open lane, hence infinitely many
independent fair bits.  Its binary law is nonatomic, and so is its
image on $[0,1]$.

If a fiber has any closed position, that position repeats along an
infinite arithmetic progression of indices.  All the corresponding
bits are forced to zero.  The set of binary sequences with zeros on
some infinite arithmetic progression is a countable union of
fair-coin null sets.  Its image has Lebesgue measure zero.  Every
fiber with a closed position is carried by this same null set.
This common support, together with non-atomicity, proves the exact
decomposition \eqref{samp:decomposition}.

To prove full support, choose any $k$ distinct successful primes
$s_0,\ldots,s_{k-1}$, and take $R>\max s_j$.  At disabled primes
$q\le R$ impose $V=0$, $U\ne0$ modulo $q$.  At successful primes
$q\le R$ impose $V\ne0$; at $s_j$ also impose $U+jV=0$.
These compatible local conditions have positive probability.
They make each of the first $k$ positions open: it has a successful
prime divisor below $R$ and no disabled divisor there.
The independent fair bits can then realize any prescribed word.
Every dyadic interval has positive mass.

Finally, impose the same conditions below $R$, without prescribing
the first $k$ positions.  Call this positive-probability event $D_R$.
For any fixed local residues on $D_R$, every closed index belongs to
the periodic set avoiding the single roots at all successful
primes $q\le R$.  This set has density
\[
 \varepsilon_R=\prod_{\substack{q\le R\\q\notin\mathcal F}}(1-1/q).
\]
With $Q_R=\prod_{q\le R}q$, at least
$(1-\varepsilon_R)n-Q_R$ of the first $n$ bits are consequently
open.  The submeasure
$\mu_R(B)=\Prob(Z_g\in B,\ D_R)$ satisfies, for every dyadic
interval $I$ of length $2^{-n}$,
\[
 \mu_R(I)\le 2^{Q_R}2^{-(1-\varepsilon_R)n}.
\]
The same bound with a different constant holds for arbitrary
intervals, by covering one interval with at most two dyadic
intervals of comparable length.  Any full-$\nu_g$-measure Borel
set has positive $\mu_R$ mass.  Covering it by intervals and summing
this bound shows that its Hausdorff dimension is at least
$1-\varepsilon_R$.  Hypothesis \eqref{samp:success-hypothesis}
gives $\varepsilon_R\to0$, proving dimension one.  A nonatomic
full-support probability has a continuous strictly increasing
distribution function; differentiation of its Lebesgue
decomposition gives the final assertion.
\end{proof}
